\begin{document}
\maketitle

{
\setcounter{tocdepth}{2}
\tableofcontents
}
\section{Zusammenfassung}\label{zusammenfassung}

Die Topological Fixed-Point Theory (TFPT) verbindet eine markierte
algebraische Ausgangsstruktur mit Darstellungen, Gitterdaten und
Kandidaten für physikalische Größen. Der Universalraum bezeichnet das
weitergehende Ziel, Zustände, erlaubte Operationen, Aufzeichnungen und
beobachtbare Entwicklung als Teile einer gemeinsamen Konstruktion zu
verstehen. Die vorliegende Arbeit synthetisiert die bis zum 15.
September 2026 verfügbaren Hauptdokumente, Fortsetzungen und
ausgewählten aktuellen Forschungsartefakte. Sie integriert insbesondere
die beiden verschiedenen v1.6.6-Linien und die neuere v1.6.7 mit ihrer
korrigierten Hamilton-Lanczos-Kette.

Der belastbare mathematische Kern umfasst eine markierte
\(D_5\)--\(A_3\)-Verklebung zu \(E_8\), eine konkrete Darstellung
desselben Gitters in einer Ordnung von \(M_2(\mathbb Q(i))\), endliche
Kontextprozesse, kohärenzverträgliche Aufzeichnungen und eine native
Paarwechselwirkung zwischen 64 Fermion- und 60 Bosonmoden. Für den
ausdrücklich festgelegten nativen Hamiltonoperator liegen Grundzustands-
und Antwortsätze vor. Die jüngste Fortsetzung bestimmt vier
Singulett-Richtungen auf der zweiten Bosonstufe, fünf korrekte
Lanczos-Glieder und zehn Momente der gefüllten Referenz. Sie verschärft
eine Energieobergrenze, ohne die gesamte Bank diagonalisiert zu haben.
Der jüngste Nachtrag identifiziert zusätzlich einen nativen
Rang-60-Kodierer, prüft dessen Grenze bei Modenverlust und klassifiziert
eine Z4-verträgliche Hamiltonfamilie bis zu einer konkreten kanonischen
Reduktion.

Der eigene Beitrag dieser Synthese besteht aus drei Verbindungen.
Erstens wird ein gemeinsamer Prozess durch den positiven Gramkern seiner
zulässigen Operationswörter rekonstruiert; daraus folgt ein präzises
Kriterium dafür, wann zwei vermeintliche Teillösungen tatsächlich
dieselbe Ausführung beschreiben. Zweitens erhält der bisher zusätzliche
Vierzustands-Transfer eine explizite, achtdimensionale Erweiterung mit
einer zweistufigen Ladungsreferenz. Sie erhält die betrachtete
Gesamt-Cartanladung exakt und reproduziert den bekannten Transferblock;
sie beweist nicht die Herkunft dieser Referenz. Drittens werden
Überlappung, dynamische Kompression und Projektorgeometrie
zusammengeführt. Ein bewegter physischer Unterraum kann geometrische
Krümmung tragen, während eine bloße Basisrotation im festen
vollständigen Raum keine erzeugt. Ein konkreter Zweikomponentenzeuge
liefert sowohl Krümmung als auch Zustandsmetrik und eine exakt bezahlte
selektive Schleifenwahrscheinlichkeit.

Diese Ergebnisse liefern einen mathematisch zusammenhängenden Ansatz für
weitere Forschung. Sie schließen die physische Auswahl von Dynamik und
Zustand, den räumlichen Ursprung, die relativistische Feldtheorie oder
die acht vollständigen TFPT-Abschlussaufgaben T1--T8 nicht. Insbesondere
folgt weder eine Theory of Everything noch ein Beweis der Riemannschen
Vermutung oder eine Lösung von P versus NP. Die offene Kernfrage lässt
sich jedoch deutlich schärfer formulieren: Welche ursprüngliche Regel
wählt einen positiven, komponierbaren Prozesskern samt lokal
zugänglichen Operationen und einer geeigneten Größenfolge aus?

\textbf{Schlüsselbegriffe:} TFPT, Universalraum, \(E_8\),
Prozessrekonstruktion, native Fockdynamik, Gramoperator, relationale
Ladung, Lanczos, geometrische Phase, physikalische Herkunft.

\begin{figure}
\centering
\pandocbounded{\includegraphics[keepaspectratio,alt={Die gemeinsame Kette und ihre Beweisgrenzen. Die drei unteren Felder unterscheiden vorliegende Modellresultate, eigene Verbindungen und die noch fehlende physische Ableitung.}]{figures/01_gemeinsame_kette.pdf}}
\caption{Die gemeinsame Kette und ihre Beweisgrenzen. Die drei unteren
Felder unterscheiden vorliegende Modellresultate, eigene Verbindungen
und die noch fehlende physische Ableitung.}
\end{figure}

\section{1. Fragestellung, Quellenstand und wissenschaftliche
Reichweite}\label{fragestellung-quellenstand-und-wissenschaftliche-reichweite}

\subsection{1.1 Was „ganzheitlich`` hier
bedeutet}\label{was-ganzheitlich-hier-bedeutet}

Eine ganzheitliche Theorie benötigt mehr als Ergebnisse mit ähnlichen
Dimensionszahlen. Sie muss zeigen, dass ihre Algebra, ihr Zustand, ihre
Dynamik und ihre Messungen zusammenpassen. Wenn ein Grundzustand in
Modell A, ein Transport in Modell B und ein chirales Feld in Modell C
berechnet werden, entsteht daraus erst dann eine gemeinsame Physik, wenn
explizite Abbildungen die jeweiligen Annahmen und beobachtbaren Prozesse
miteinander verbinden.

Daher lautet die Leitfrage dieses Papers:

Welche einheitlichen Ausgangsdaten bestimmen
\((\mathcal A,\alpha_t,\omega,\mathfrak I)\) und deren physikalische
Interpretation?

Hier ist \(\mathcal A\) eine Algebra von Observablen, \(\alpha_t\) ihre
Entwicklung, \(\omega\) ein Zustandsfunktional und \(\mathfrak I\) die
Familie tatsächlich zugelassener Instrumente. Ein Instrument umfasst die
möglichen Messausgänge samt Zustandsänderung und Wahrscheinlichkeit. Die
Notation ist eine Spezifikation; sie wählt diese Daten noch nicht aus.

Anschaulich beschreibt die bisherige TFPT-Grammatik, welche Bauteile
zusammenpassen. Eine vollständige Ausführung muss zusätzlich erklären,
welche Handgriffe vorkommen, was anfänglich vorliegt und welche Spuren
die Handgriffe hinterlassen. Raum wäre dann eine besondere Struktur
unabhängiger Zugriffsmöglichkeiten und endlicher Signalwege. Zeit wäre
mehr als die Ordnung einer zyklischen Matrix: Sie müsste als gemeinsame
Entwicklung mit lesbaren Korrelationen auftreten.

\subsection{1.2 Quellenhierarchie und
Stichtag}\label{quellenhierarchie-und-stichtag}

Der Ausgangskorpus umfasst 74 eingefrorene Dokumentdateien
einschließlich des nachgereichten Texts. Darunter befinden sich elf PDFs
mit zusammen 830 Seiten; ihre Texte wurden extrahiert. Die zentralen und
neu geänderten Abschnitte wurden gezielt gelesen, repräsentative
Ausgangsseiten visuell kontrolliert. Die Zahl der inventarisierten
Seiten bedeutet ausdrücklich keinen vollständigen erneuten Beweisaudit
jeder historischen Aussage.

Die Hauptquellen sind:

\begin{itemize}
\tightlist
\item
  \textbf{S046:} Hauptdokument v1.6 vom 15. September, 168 PDF-Seiten;
  breite historische und physikalische Grundlage.
\item
  \textbf{S033:} Clock/Gemeinsame-Quelle-Konsolidierung v1.6.6; innere
  Spin-Identifikation, Operationsverträge, Transferzeuge und
  Feldkorrekturen.
\item
  \textbf{S001:} andere v1.6.6-Linie, „Fundamentale Fortsetzung``;
  Präparationslücke, quantitative Restantwort, Trennschnitte und
  Feldstabilität.
\item
  \textbf{S003:} fundamentale Reduktion; symmetrische Gesamtantwort,
  Multiplizitätsräume und operationaler Wirkungstest.
\item
  \textbf{S068/S070:} Ergebnisbericht und vollständige Konsolidierung
  v1.6.7; korrekte H-Kette, Singulettprojektionen, bilineare Zerlegung
  und Unterbestimmtheit.
\item
  \textbf{S074:} jüngster Nachtrag derselben v1.6.7 mit nativer
  Kodierung, Schattenrekonstruktion, Interventionsgegenbeispiel und
  Z4-/Bogoliubov-Analyse.
\item
  \textbf{S071/S072:} zentrale offene Probleme und Compiler-Synthese des
  Repositorys.
\end{itemize}

S001 und S033 sind verschiedene Dokumentlinien trotz gleicher
Versionsnummer. Inhalt, Tensorpin und Beweisvoraussetzungen entscheiden
über ihre Kombination. S070 bewahrt große Teile älterer Fassungen als
historische Anhänge; das bloße Wiederauftauchen einer Aussage ist kein
neuer Beweis. In der vorliegenden Synthese stehen jeweils die neuesten
kompatiblen Aussagen im Haupttext.

Das Repository hatte bei der Bestandsaufnahme den Commit
\texttt{66b91e40e245} (vollständiger Hash im Prüfpaket), aber viele
neuere, nicht eingecheckte Forschungsdateien. Deshalb bezeichnet der
Commit allein den hier ausgewerteten Stand nicht. Das Quellenmanifest
hält die tatsächlich verwendeten Dateien mit SHA-256 fest. Aktive fremde
Forschungsarbeiten wurden nicht verändert. Die Hauptsynthese
berücksichtigt auch den während der Arbeit hinzugekommenen Nachtrag
S074. Die früher eingefrorenen Fassungen bleiben zum Versionsvergleich
erhalten; der letzte inhaltliche Nachtrag wurde am 15.09.2026 gegen
12:30 Uhr MESZ gesichert.

\subsection{1.3 Beweisstatus}\label{beweisstatus}

\textbf{Gesetzt} sind Modellannahmen und Ressourcenzugänge.
\textbf{Exakt} ist eine Identität oder ein mathematischer Schluss unter
genannten Voraussetzungen. \textbf{Numerisch} bezeichnet eine
Approximation ohne rigoroses Intervall. \textbf{Bedingt} bedeutet, dass
ein Resultat zusätzliche, noch nicht hergeleitete Eingaben benötigt.
\textbf{Offen} kennzeichnet fehlende Beweise oder physikalische
Identifikationen.

„Exakt`` und „bedingt`` können zugleich gelten. Beispielsweise ist die
Ladungserhaltung der neuen Referenzkonstruktion exakt; dass TFPT diese
Referenz erzeugt, bleibt offen. Ein bewiesener Satz über ein
festgelegtes Modell ist außerdem noch keine empirische Bestätigung
dieses Modells.

Die eigenen Rechnungen umfassen 57 explizite exakte Prüfbedingungen,
darunter native Tensoridentitäten, eine vollständige
Graphcharakteristik, rationale Ritz-Einschließung, symbolische
Referenzkopplung und Projektorgeometrie. Normaler und optimierter Lauf
stimmen bytegenau überein. Zusätzlich wurde der 20 Bedingungen
umfassende H-Ketten-Prüfer aus v1.6.7 in einer isolierten Kopie
ausgeführt. Der große historische Grundzustandslauf und die vollständige
Vier-Boson-Enumeration wurden hier nicht erneut gerechnet. Ihre
Voraussetzungen bleiben entsprechend sichtbar.

\section{2. Die algebraische Grundlage: was TFPT tatsächlich
zusammenbindet}\label{die-algebraische-grundlage-was-tfpt-tatsuxe4chlich-zusammenbindet}

\subsection{2.1 Ausgangsstruktur und
Einheiten}\label{ausgangsstruktur-und-einheiten}

Die historische Kurzfassung organisiert sich um

\[c_3=\frac1{8\pi},\qquad g_{\mathrm{car}}=5.\]

Der zweite Ausdruck ist eine Trägerzahl, nicht die Kopplung \(g\) des
späteren Hamiltonoperators. Die erste Eingabe enthält mehr als eine
Zahl: orientierte Rand- oder Nahtstruktur, einen geeigneten positiven
Kern und eine Normierung. Die zweite benötigt einen konkreten komplexen
Träger und Markierungen. Diese Daten dürfen nicht verschwinden, wenn man
die Theorie verbal auf „zwei Zahlen`` verkürzt. {[}S046, Kap. 3; S072,
§2{]}

Der fünfteilige Träger liefert

\[S^+=\Lambda^{\mathrm{even}}\mathbb C^5,
\qquad \dim S^+=1+10+5=16.\]

Das ist die chirale Spinordarstellung der inneren Gruppe
\(\mathrm{Spin}(10)\). Daraus folgt noch nicht, dass die physikalische
Raumzeit zehn- oder fünfdimensional wäre. Ebenso ist der arithmetische
Familienzähler \((16-1)/5=3\) kein Beweis dreier leichter chiraler
Familien.

Eine dimensionslose Theorie kann Verhältnisse vorhersagen. Die Angabe
einer Energie in GeV benötigt zusätzlich einen Bezug zu einer
dimensional kalibrierten Größe. Einheitenwahl und Vorhersage
dimensionsloser Verhältnisse sind getrennte Fragen; die Verwendung einer
eingesetzten Planckskala muss im physikalischen Anschluss dokumentiert
bleiben.

\subsection{\texorpdfstring{2.2 Von \(D_5\oplus A_3\) zu
\(E_8\)}{2.2 Von D\_5\textbackslash oplus A\_3 zu E\_8}}\label{von-d_5oplus-a_3-zu-e_8}

Beide Wurzelgitter haben Determinante vier. Eine geeignete diagonale
Verklebung ihrer Diskriminantklassen erzeugt einen geraden Überverband
vom Index vier. Seine Determinante ist

\[\det L=\frac{\det D_5\det A_3}{4^2}=1.\]

Mit Positivität und Rang acht erhält man das gerade unimodulare
\(E_8\)-Gitter. Die Aussage benötigt die Verklebungsform; die Rangsumme
\(5+3=8\) allein reicht nicht. Die 240 Wurzeln sind 112 ganzzahlige
Wurzeln mit zwei Einträgen \(\pm1\) und 128 Halbvektoren mit gerader
Minusparität. {[}S046, Kap. 3{]}

Die komplexe Lie-Algebra zerfällt als

\[
\mathfrak e_8=(45,1)\oplus(1,15)\oplus(16,4)
\oplus(\overline{16},\overline4)\oplus(10,6).
\]

Die Dimensionen summieren sich zu 248. Insbesondere treten dieselben 64
und 60 auf, die später den nativen Paarvertrag tragen. Diese
Übereinstimmung ist strukturell relevant, weil ein konkreter Intertwiner
die Paarung realisiert. Sie ersetzt dennoch keine Auswahl von Energie,
Zustand oder räumlicher Ausführung.

\subsection{2.3 Der kleine Matrixkern}\label{der-kleine-matrixkern}

Im gemeinsamen Matrixkörper \(M_2(\mathbb Q(i))\) seien

\[
a=iI_2,\quad
u_1=\begin{pmatrix}i&0\\0&-i\end{pmatrix},\quad
u_2=\begin{pmatrix}0&1\\-1&0\end{pmatrix},\quad
u_3=u_1u_2,
\]

Die Familie wird durch

\[w=\frac12(I+u_1+u_2+u_3),\qquad w^3=-I\]

ergänzt. Mit \(D=\mathbb Z[i]\) und

\[P=\begin{pmatrix}1&(1+i)^{-1}\\0&(1+i)^{-1}\end{pmatrix}\]

liegt eine maximale Ordnung \(\mathcal M=PM_2(D)P^{-1}\) vor. Unter der
reellen Form

\[\langle A,B\rangle=\operatorname{Re}\operatorname{Tr}(AB^\dagger)\]

hat ihre ganzzahlige Acht-Basis positive führende Hauptminoren
\(2,4,4,4,4,4,2,1\). Die Form ist gerade und unimodular: wiederum
\(E_8\). Die Quellen enthalten außerdem eine markierungstreue Abbildung
zum tatsächlich verwendeten Compiler-Gitter. {[}S046, Kap. 34{]}

Das ist stärker als eine zufällige Zahlengleichheit. Gleichwohl sind nur
96 der 240 Wurzelmatrizen unitär; 144 haben Rang eins. Eine Wurzel ist
somit nicht automatisch eine ausführbare reversible Zeitoperation.
Verschiedene Ordnungen können dieselbe hermitesche Kegelansicht haben
und sich in ihren nicht-hermiteschen Operationen unterscheiden. Der
positive Kegel wählt den vollständigen Operationssatz nicht allein aus.

\subsection{2.4 Der Clock-Anschluss}\label{der-clock-anschluss}

Die neuere native Clock-Konstruktion verwendet eine Permutation auf fünf
Hilfsachsen und ihre vorzeichenrichtige Exteriorhebung. Mit zehn
Majorana-Matrizen \(\gamma_j\), der Parität \(\Pi\) und

\[R_{ab}=\tfrac12(\gamma_a-\gamma_b)(\gamma_{a+5}-\gamma_{b+5})\]

wird explizit

\[S=\gamma_0\cdots\gamma_9R_{01}R_{02}R_{34}
=\Pi\Gamma(p)\in\mathrm{Spin}(10).\]

Die Wirkung auf gerade Spinoren und Vektoren reproduziert die beiden
nativen Darstellungen und erhält den Tensor. Dieser Clock hat Ordnung
sechs. {[}S033, §3{]}

Frühere Ordnung-3-Familienoperationen und Ordnung-12-Lifts gehören zu
ihren jeweils bezeichneten Darstellungen. Matrixordnung, projektive
Ordnung und Ordnung der Konjugationswirkung sind getrennt zu prüfen.
Ihre Zahlen werden hier nicht zu einer einzigen universellen
Periodizität zusammengezogen. Der neue Satz verbindet einen konkreten
Clock mit der inneren Symmetrie; er liefert weder kontinuierliche
Symmetriekontrollen noch eine Zeiteinheit.

\section{3. Universalraum als Prozess: Zustand, Gedächtnis und
Schatten}\label{universalraum-als-prozess-zustand-geduxe4chtnis-und-schatten}

\subsection{3.1 Das endliche
Kontextmodell}\label{das-endliche-kontextmodell}

Im markierten komplexen Viererträger liegen 60 Strahlen in 15
orthogonalen Viererkontexten vor. Bei der festgelegten Inzidenzregel
\(B\) und gleichmäßiger Wahl der sieben zulässigen Nachfolgekontexte
gilt

\[
T_{(D,t),(C,s)}=\frac{B_{DC}}7
\operatorname{Tr}(\Pi_{D,t}\Pi_{C,s}).
\]

Jeder Ausgang besitzt einen Übergang mit Gewicht \(1/7\) und zwölf mit
Gewicht \(1/14\). Die genaue Faktorisierung führt zu Rang 30 und dem
Spektrum

\[\{1^{[1]},(2/7)^{[9]},(-2/7)^{[5]},(3/7)^{[15]},0^{[30]}\}.\]

Die Born-Regel, die zulässigen Kontexte und deren Gewichtung sind dabei
Eingaben. Der stationäre Systemzustand \(I_4/4\) ist ein Fixpunkt dieses
Messprozesses; er ist nicht bereits das physikalische Vakuum. {[}S046,
Kap. 4{]}

\subsection{3.2 Warum eine einzige Dichtematrix nicht immer
genügt}\label{warum-eine-einzige-dichtematrix-nicht-immer-genuxfcgt}

Für klassische Kontextinformation und ein Quantensystem kann man einen
Zustand durch positive Blöcke \(X_C\) beschreiben. Die feste Ausführung
lautet

\[\Phi(X)_D=\sum_C K_{DC}\Delta_D(X_C).\]

Die Quellen zeigen \(\operatorname{rank}\Phi=60\) und
\(\operatorname{rank}\Phi^2=30\). Für eine bestimmte passive
Marginalauslesung schließen sogar 45 lineare Größen. Dagegen können
selektive Eingriffe oder wiederverwendete kohärente Register mehr
Information zugänglich machen. Die minimale Beschreibung hängt somit vom
tatsächlichen Operations- und Auslesevertrag ab. {[}S046, Kap. 5{]}

Ein System kann nach einer Vormessung dieselbe reduzierte Dichtematrix
haben, obwohl sein gemeinsamer Zustand mit einem behaltenen Register
eine andere spätere Antwort liefert als mit einem frischen Register. Das
ist kein Widerspruch: Die gemeinsame Korrelation fehlt in der
reduzierten Beschreibung. Die etablierte Theorie von Quantennetzwerken
und Prozesstensoren beschreibt genau solche Mehrzeitverträge. {[}E1,
E2{]}

